\chapter{HR-MMoH: 多模态大语言模型在异构 GPU 集群上的重计算策略}

第二章分析了多模态大语言模型（MLLMs）在模块异构与分阶段冻结训练下呈现的计算与显存需求差异，第三章在此基础上通过 MMoH 框架解决了异构集群上的划分与调度问题。
在此背景下，本章进一步聚焦显存维度的瓶颈，提出 HR-MMoH（Heterogeneous Recomputation for training multimodal large language models on heterogeneous GPU clusters），
以异构感知的重计算策略补充第三章的工作。

\section{引言}
\label{sec:HR-MMoH-intro}

第三章介绍的 MMoH 框架通过需求驱动的设备拓扑、复合粒度划分与提前终止流水线调度，在异构 GPU 集群上显著提升了训练 MLLMs 的吞吐量。
然而，由于异构 GPU 集群中设备间显存与算力差异显著，在显存受限的节点上，仅靠流水线模型划分仍可能无法容纳较大批量的数据或者较长的序列输入，训练过程仍面临显存瓶颈。
\emph{重计算}（Recomputation）又称\emph{梯度检查点}（Gradient Checkpointing）\cite{chen2016checkpoint}，
通过在前向传播时少存部分中间激活值或者不存任何中间激活值、在反向传播时重新进行一次前向传播，把在前向传播时丢弃的中间激活值重新计算出来，
以算力换显存，在保持数学等价性的前提下支持更大批量、更长序列或更大参数的模型训练，是目前学术界和工业界广泛采用的缓解显存压力的技术。

现有主流框架（如 Megatron-LM）中的重计算策略多面向\emph{同构}集群设计：在同一训练任务中，
对所有流水级或所有设备采用统一的\emph{重计算粒度}（如整层重计算 Full Recomputation、子模块级重计算 Selective Recomputation）
与统一的\emph{层级划分方式}（如均匀分块 uniform、仅前若干层做检查点 block），
或对全体层采用相同的\emph{重计算子模块集合}（如仅 core\_attn、或 core\_attn+mlp+layernorm）。

然而，在异构集群上，各设备显存与算力差异显著：若在显存紧张的设备上采用过粗的重计算粒度或过少的重计算模块，
仍易发生 OOM，导致训练无法进行；
若在显存充裕的设备上采用过细的粒度或过多的重计算子模块，则会带来不必要的重计算开销，
延长迭代时间并造成算力浪费。
因此，\emph{为了同构集群设计的的检查点策略难以在异构环境下同时兼顾显存利用率与计算效率}。

针对上述问题，本章介绍 HR-MMoH：
在 MMoH 的异构流水线划分与调度基础上，面向\emph{异构设备}的显存与算力差异，
设计\emph{异构感知的重计算}策略，使各流水级或各设备可根据自身显存与算力条件
采用差异化的重计算粒度与子模块重计算选择，在降低 OOM 风险的同时尽量减少多余重计算、平衡各流水级计算负载，
从而与第三章的划分与调度形成“划分—调度—显存优化”的联合优化。
\ref{sec:HR-MMoH-method}节形式化描述异构重计算策略的设计与实现思路，
\ref{sec:HR-MMoH-exp}节给出实验设置、结果与分析，\ref{sec:HR-MMoH-summary}节总结本章内容。

\section{异构重计算}
\label{sec:HR-MMoH-method}

\subsection{现有的重计算策略与定量分析}
\label{sec:megatron-recompute}

作为当前大模型分布式训练的主流框架，NVIDIA 推出的 Megatron-LM\cite{megatron2020,megatron2021,megatron2023} 大模型并行训练框架集中定义了与重计算相关的配置，
开发者通过命令行参数传入重计算配置，例如 \texttt{--recompute-granularity}、\texttt{--recompute-method}、\allowbreak \texttt{--recompute-num-layers}、\texttt{--recompute-modules}。

配置参数方面，表~\ref{tab:chap4-megatron-params} 给出 Megatron 中与重计算相关的主要参数及其含义。
其中 \texttt{recompute\_granularity} 决定整层（full）或子模块级（selective）重计算；
full 时须同时指定 \texttt{recompute\_method} 与 \texttt{recompute\_num\_layers}，selective 时二者须为 \texttt{None}，仅由 \texttt{recompute\_modules} 指定要重算的子模块。
\texttt{distribute\_saved\_activations} 为 \texttt{True} 时，可将 checkpoint 保存的激活在张量并行组内分片存储以进一步省显存，但仅对 full 粒度且未启用序列并行的情形有效。

\begin{table}[htbp]
	\centering
	\caption{Megatron-LM 重计算相关配置参数}
	\label{tab:chap4-megatron-params}
	\begin{tabular}{@{}p{3.5cm}p{2.8cm}p{\dimexpr\textwidth-3.5cm-2.8cm-4\tabcolsep\relax}@{}}
		\toprule
		参数 & 类型/取值 & 说明 \\
		\midrule
		\texttt{recompute\_}\\\texttt{granularity} & \texttt{None} $\mid$ \texttt{full} $\mid$ \texttt{selective} & 重计算粒度。\texttt{None} 表示不重计算；\texttt{full} 对整层做重计算；\texttt{selective} 仅对 \texttt{recompute\_modules} 中列出的子模块重算。 \\
		\texttt{recompute\_method} & \texttt{None} $\mid$ \texttt{uniform} $\mid$ \texttt{block} & 仅用于 \textbf{full} 粒度：如何划分“哪些层”做 checkpoint。\texttt{uniform} 为均匀分块，\texttt{block} 为仅对每流水级前若干层做 checkpoint。\textbf{selective 时必须为 None}。 \\
		\texttt{recompute\_num\_}\\\texttt{layers} & \texttt{None} $\mid$ 正整数 & 仅用于 \textbf{full} 粒度。\texttt{uniform} 时表示每个 chunk 的层数；\texttt{block} 时表示每个 stage 内做重计算的层数。\textbf{selective 时必须为 None}。 \\
		\texttt{distribute\_saved\_}\\\texttt{activations} & 布尔 & 为 \texttt{True} 时，将 checkpoint 保存的激活在张量并行组内分片存储；仅 \textbf{full} 粒度且非序列并行时有效。 \\
		\texttt{recompute\_modules} & 字符串列表 & 仅用于 \textbf{selective}。要重算的子模块名列表（见下文）。默认 \texttt{["core\_attn"]}。 \\
		\bottomrule
	\end{tabular}
\end{table}

\textbf{整层重计算（full）}：以 Transformer 层为统一粒度做检查点；前向只保存每个“块”的输入，反向时重新执行该段前向再反传。
在 full 粒度下，\texttt{recompute\_method} 与 \texttt{recompute\_num\_layers} 共同决定流水级上多少个模型层只需要保存一个中间激活值。
\texttt{recompute\_method} 具体可以取一下以下两种值：
\begin{itemize}
	\item \textbf{uniform}（均匀分块）：将当前流水级内的所有层均匀分成多个 chunk，每个 chunk 包含 \texttt{recompute\_num\_layers} 层；
  每个 chunk 的输入被 checkpoint，chunk 内前向一次执行，反向时整个 chunk 先重计算再反传。
  chunk 数少则显存更省，但单次重算层数较多。
	\item \textbf{block}（仅前 N 层）：在当前流水级内，只对前 \texttt{recompute\_num\_layers} 层逐层做重计算；后续层不需要做重计算，正常存储中间激活。
  在显存允许时可减少重计算量，更充分利用显存。
\end{itemize}

\textbf{选择重计算（selective）}：在每一层内仅对 \texttt{recompute\_modules} 中列出的子模块做重计算，其余部分正常存激活。
此时 \texttt{recompute\_method} 与 \texttt{recompute\_num\_layers} 均须为设置为 \texttt{None}，且 selective 始终作用于所有流水级上的所有模型层。
模型每个子模块在各自的前向传播中根据用户指定的配置项\texttt{recompute\_modules} 决定舍弃哪些子模块的中间激活值，在反向传播时哪些模块的激活值需要重新计算。

选择性子模块方面，\texttt{recompute\_modules} 指定在 selective 粒度下要对哪些子模块做重计算。
不同子模块采用的重计算机制不同：标准 checkpoint 保存输入、反向时重新前向再反传；
Output-discarding checkpoint 在前向完成后丢弃该子模块输出的显存，在反向时通过梯度钩子函数再算一遍并写回，可进一步节省显存。
若未指定 \texttt{recompute\_modules}，则默认为 \texttt{["core\_attn"]}。
稠密模型下常用的“最省显存”组合为 \texttt{[core\_attn, mlp, layernorm]}。
表~\ref{tab:chap4-megatron-modules} 列出目前训练框架使用选择重计算时支持的常用选项、含义及 checkpoint 类型。

\begin{table}[htbp]
	\centering
	\caption{Megatron selective 粒度下可配置的子模块}
	\label{tab:chap4-megatron-modules}
	\renewcommand{\arraystretch}{1.2}
	\begin{tabular}{lll}
		\toprule
		模块名 & 含义 & Checkpoint 类型 \\
		\midrule
		\texttt{core\_attn} & 核心注意力（QKV$\to$Attention$\to$输出） & 标准 \\
		\texttt{mlp} & 稠密 MLP（非 MoE） & 标准 \\
		\texttt{layernorm} & 输入 LayerNorm 与 MLP 前 LayerNorm & Output-discarding \\
		\texttt{moe} & 整块 MoE 层（router + experts + combine） & 标准 \\
		\texttt{shared\_experts} & MoE 的共享专家 & 标准 \\
		\texttt{moe\_act} & MoE 专家内部激活（如 SiLU）与专家权重 & Output-discarding \\
		\texttt{mla\_up\_proj} & MLA 的 QKV 上投影与 RoPE & Output-discarding \\
		\bottomrule
	\end{tabular}
	\renewcommand{\arraystretch}{1.0}
\end{table}


上述策略在实现上对同一训练任务内的所有流水级、所有 GPU 计算设备采用统一的重计算粒度、重计算方法、重计算层数以及重计算模块，即面向同构集群时使用了的“全局一致”的重计算策略。
但是在异构集群上，显存紧张流水级若与显存充裕流水级使用相同的重计算配置，易出现前者仍 OOM 或后者重计算过多的问题，
而且目前的重计算技术给所有流水级增加了同样大小的计算负载，这使得异构 GPU 集群中性能较差的 GPU 进一步成为计算瓶颈，
为下文针对各流水级/设备做差异化配置的异构重计算提供了动机。

\paragraph{重计算技术的定量分析。}
为定量比较不同重计算策略对显存占用的影响，本文在给定输入形状与数值格式下，
对单个 Transformer 层在前向传播中需为反向传播保留的中间激活所需要的显存占用量进行详细推导。
假设层输入张量 X 形状为 $[B, S, H]$，其中 $B$ 为 batch size、$S$ 为输入序列长度、$H$ 为隐藏维度；
中间激活使用训练常用的 BF16 格式，每元素占 2 个字节。Transformer 层采用常见的 Pre-LayerNorm 的 decoder-only 单层结构：
层输入 X 经 LayerNorm 后进入注意力子层，
残差连接后再经 LayerNorm 进入 MLP 子层（线性层 $\to$ GELU激活层 $\to$ 线性层 + Dropout 层），再次残差连接得到层输出。
记注意力头数为 $A$，MLP 层的中间维度设置为常用的设置，即隐藏维度的 4 倍，记为 $4H$。

为避免后续公式中大小写符号混用带来的歧义，表~\ref{tab:chap4-symbols} 汇总本小节使用的主要记号。

\begin{table}[htbp]
  \centering
  \caption{4.2.2 节主要符号说明}
  \label{tab:chap4-symbols}
  \small
  \renewcommand{\arraystretch}{1.15}
  \begin{tabular}{@{}p{2.6cm}p{10cm}@{}}
    \toprule
    符号 & 含义 \\
    \midrule
    $B$ & batch size \\
    $S$ & 输入序列长度 \\
    $H$ & 模型隐藏维度 \\
    $A$ & 注意力头数。 \\
    X & 输入张量，形状为 $[B,S,H]$。 \\
    $\mathbf{h}_1$ & 注意力子层与残差连接后的中间张量，形状为 $[B,S,H]$。 \\
    $N_{\mathrm{attn}}$ & 单个注意力子层为反向传播保留的中间激活总量。 \\
    $N_{\mathrm{mlp}}$ & 单个 MLP 子层为反向传播保留的中间激活总量。 \\
    $N_{\mathrm{ln}}$ & 单个 LayerNorm 子层为反向传播保留的中间激活总量。 \\
    $N_{\mathrm{None}}$ & 不采用重计算时，单层 Transformer 需保留的中间激活总量。 \\
    $N_{\mathrm{selective}}^{\mathrm{core\_attn}}$ & 仅对 \texttt{core\_attn} 使用 selective 重计算时，单层需保留的中间激活总量。 \\
    $N_{\mathrm{selective}}^{\mathrm{core\_attn+mlp}}$ & 对 \texttt{core\_attn} 与 \texttt{mlp} 使用 selective 重计算时，单层需保留的中间激活总量。 \\
    $N_{\mathrm{selective}}^{\mathrm{core\_attn+mlp+ln}}$ & 对 \texttt{core\_attn}、\texttt{mlp} 与 \texttt{layernorm} 使用 selective 重计算时，单层需保留的中间激活总量。 \\
    $N_{\mathrm{full}}$ & 采用 full 重计算时，单层需保留的中间激活总量。 \\
    \bottomrule
  \end{tabular}
  \renewcommand{\arraystretch}{1.0}
\end{table}

\textbf{（1）不重计算（\texttt{None}）。} 仅计反向传播所需中间激活， 使用训练时常用的 BF16 格式，所以每个激活参数大小为 2 字节。

\paragraph{注意力子层}：
Q/K/V 共享输入为 $2 B S H$，计算得到的 Q/K/V 大小均为 $2 B S H$，共计 $6 B S H$，
计算 $QK^{T}$ 后进行 softmax 的中间结果为 $2B A S^2$，Dropout 的 mask 为 $B A S^2$，
Dropout 的输出为 $2 B A S^2$，最后与 V 相乘的中间结果为 $2 B A S$，
最后的 Dropout 的 mask 矩阵为 $B S H$，
合计
\begin{equation}
  N_{\mathrm{attn}} = 11 B S H + 5 B A S^2.
\end{equation}
其中 $5 B A S^2$ 来自 $QK^{\top}$、softmax 及 dropout等，随序列长度平方增长，成为长序列输入时中间激活的主要来源。
图~\ref{fig:chap4-recomp-attn} 详细分析了注意力子层在 \texttt{None} 情形下需保存的中间激活及其大小。

\begin{figure}[htbp]
  \centering
  \includegraphics[width=0.85\linewidth]{recomputation_attention.pdf}
  \caption{注意力子层中间激活示意}
  \label{fig:chap4-recomp-attn}
\end{figure}

\paragraph{MLP 子层}：
中间维度 $4H$，第一个线性层的输入为 $2 B S H$，需要保存的输出为 $8 B S H $，经过中间激活后第二个线性层的输入为 $8 B S H$，
最后的 Dropout 层的 mask 矩阵为 $B S H$，合计
\begin{equation}
  N_{\mathrm{mlp}} = 19 B S H.
\end{equation}
图~\ref{fig:chap4-recomp-mlp} 给出了 MLP 子层在 \texttt{None} 情形下需保存的中间激活及其大小。

\begin{figure}[htbp]
  \centering
  \includegraphics[width=0.85\linewidth]{recomputation_mlp.pdf}
  \caption{MLP 子层中间激活示意}
  \label{fig:chap4-recomp-mlp}
\end{figure}

\emph{两个 LayerNorm}（注意力前、MLP 前）各保存输入 $2 B S H$，合计
\begin{equation}
  N_{\mathrm{ln}} = 4 B S H.
\end{equation}

于是，不采用重计算时单层需保存的中间激活总量为
\begin{equation}
  N_{\mathrm{None}} = N_{\mathrm{attn}} + N_{\mathrm{mlp}} + N_{\mathrm{ln}} = 34 B S H + 5 B A S^2.
\end{equation}
其中 $5 B A S^2$ 项随序列长度平方增长，可能成为长序列场景下中间激活显存的主要来源。

\textbf{（2）选择重计算（\texttt{selective}）。} 仅对指定子模块做重计算：前向只保存该子模块的\emph{输入}，反向传播时在子模块内部重新前向传播再反传，
其它未开启重计算的子模块仍按 \texttt{None} 情形完整保存中间激活。
以仅对 \texttt{core\_attn} 开启选择重计算为例，注意力子层在 \texttt{None} 情形下需要 $N_{\mathrm{attn}} = 11 B S H + 5 B A S^2$ 大小的显存空间来存储中间激活值；
对 \texttt{core\_attn} 做选择重计算之后，Attention 块内部的中间激活（$QK^{\top}$、softmax 及其 dropout mask 矩阵等）都不再常驻显存，
只保留该子模块的输入张量（层输入和 LN 输出）以便在反向传播重新计算注意力块本身；
MLP 子层与两个 LayerNorm 的激活仍按照 $N_{\mathrm{mlp}}$ 与 $N_{\mathrm{ln}}$ 的代价保存，因此整层的线性部分从 $34 B S H + 5 B A S^2$ 约减少为 $23 B S H$。

若在此基础上进一步对 \texttt{mlp} 开启选择重计算，则 MLP 内部的 $19 B S H$ 激活也改为在反向传播时重算，只需要让输入的张量常驻显存，此时需要保存的中间激活总量进一步降低为 $6 B S H$。
再将 \texttt{layernorm} 纳入 \texttt{recompute\_modules} 时，其输出在前向结束后同样可以丢弃，仅保留必要的输入，需要保存的中间激活总量进一步降低为 $4 B S H$。
为便于与 \texttt{None} 情形直观对比，本文用
\begin{align}
  N_{\mathrm{selective}}^{\mathrm{core\_attn}} &= 23 B S H, \\
  N_{\mathrm{selective}}^{\mathrm{core\_attn+mlp}} &= 6 B S H, \\
  N_{\mathrm{selective}}^{\mathrm{core\_attn+mlp+ln}} &= 4 B S H.
\end{align}
来表示选择不同重计算模块下需要存储的中间激活占据的显存总量。
与 \texttt{None} 相比，选择重计算策略去掉了注意力主体与 MLP 主体中的大部分常驻中间激活，仅保留子模块边界的输入张量，显存开销的孔家复杂度由 $O(B S H + B A S^2)$ 降为 $O(B S H)$。

\textbf{（3）整层重计算（\texttt{full}）。} 以整层为重计算单元：前向只保存层输入 X，反向时整层重新前向传播再反向传播计算梯度。
故
\begin{equation}
  N_{\mathrm{full}} =2 B S H.
\end{equation}
这种情况下的显存占用最小，但反向时需重算整层前向，计算开销最大。

表~\ref{tab:chap4-activation} 汇总了不采用重计算、选择重计算与整层重计算时单层需保存的中间激活总量。

\begin{table}[htbp]
  \centering
  \caption{单层 Transformer 中间激活量}
  \label{tab:chap4-activation}
  \small
  \renewcommand{\arraystretch}{1.5}
  \setlength{\tabcolsep}{4pt}
  \begin{tabular}{@{}p{2.5cm}p{4.2cm}p{5.8cm}@{}}
    \toprule
    策略 & 需保存的激活（元素个数） & 说明 \\
    \midrule
    \texttt{None} & $34 B S H + 5 B A S^2$ & 整层所有中间激活均常驻 \\
    \hdashline
    \texttt{selective} [\texttt{core\_attn}] & $23 B S H$ & 仅注意力内部重算，MLP 与 LN 仍常驻 \\
    \hdashline
    \texttt{selective} [\texttt{core\_attn}, \texttt{mlp}] & $6 B S H$ & 注意力与 MLP 内部重算，仅边界与 LN 常驻 \\
    \hdashline
    \texttt{selective} [\texttt{core\_attn}, \texttt{mlp}, \texttt{layernorm}] & $4 B S H$ & 三者均重算，仅两个子模块边界的输入常驻 \\
    \hdashline
    \texttt{full} & $2 B S H$ & 仅保存层输入，反向时整层重算 \\
    \bottomrule
  \end{tabular}\\
  \renewcommand{\arraystretch}{1.0}
  \vspace{2pt}
\end{table}

\subsection{异构重计算策略设计}
\label{sec:HR-MMoH-design}

在 MMoH 已完成的流水线划分与设备映射基础上，HR-MMoH 的异构重计算核心思想是：\emph{为每个流水级（或每台计算设备）根据其显存上限以及算力，独立配置重计算粒度与子模块集合}，
使显存压力大的流水级采用更激进的检查点以降低 OOM风险，显存压力小的流水级采用较保守的配置以减少不必要的重计算开销，从而在满足显存约束的前提下最小化迭代时间并最大化 GPU 利用率。


\paragraph{输入与约束。}
在 MMoH 规划器已完成模型划分与设备映射的前提下，记流水线并行度为 $S$，流水级索引为 $s \in \{1,\ldots,S\}$，
流水级 $s$ 映射至设备 $d_s$。记设备 $d_s$ 的显存容量上界为 $M_s$，相对计算能力系数为 $C_s$（可由该设备 FLOPs 或实测单流水级迭代时间经归一化得到）。
流水级 $s$ 由若干模型块构成；其在前向与反向传播阶段的激活显存峰值及重计算开销由该流水级所采用的检查点配置 $\theta_s$ 唯一确定。
配置 $\theta_s$ 的取值为：\texttt{recompute\_granularity} $\in \{\texttt{None}, \texttt{full}, \texttt{selective}\}$；
当为 \texttt{full} 时还需指定 \texttt{recompute\_method} 与 \texttt{recompute\_num\_layers}，
当为 \texttt{selective} 时则指定 \texttt{recompute\_modules} 子模块列表。形式化地，约束条件为：$\forall s$，
在配置 $\theta_s$ 下流水级 $s$ 的激活显存峰值不超过 $M_s$（可保留安全余量）；优化目标为最小化单次训练迭代时间，
即 $\min \max_s T_s$，其中 $T_s$ 表示流水级 $s$ 在配置 $\theta_s$ 下前向、重计算与反向阶段耗时之和（含流水级间通信）。

\paragraph{问题求解。}
本文采用基于剖析（profiling）的配置生成策略。对每一流水级 $s$，在与其对应的代表设备上预先测量不同 $\theta_s$ 下的激活显存占用与单流水级执行时间 $T_s$；在满足显存约束的可行配置集中，为流水级 $s$ 选取使 $T_s$ 较小且与相邻流水级时间较为均衡的配置，以降低瓶颈流水级对整体迭代时间的主导程度。策略可归纳为两类情形：（1）\emph{显存紧张}的流水级：优先采用 \texttt{full} + \texttt{uniform} 并取较小的 \texttt{recompute\_num\_layers}，或 \texttt{selective} 并采用较完整的 \texttt{recompute\_modules}（如 \texttt{[core\_attn, mlp, layernorm]}），以降低激活显存峰值、避免 OOM；（2）\emph{显存充裕而算力或通信紧张}的流水级：采用 \texttt{None}，或 \texttt{selective} 且仅 \texttt{[core\_attn]}，或 \texttt{full} + \texttt{block} 且较小的 \texttt{recompute\_num\_layers}，以减少不必要的重计算与同步开销。实现上，在 MMoH 规划器输出流水级划分与设备映射之后，由 HR-MMoH 配置模块为各流水级生成对应的 \texttt{TransformerConfig} 参数（或等价的环境变量/配置文件），各训练进程（rank）在初始化时加载其所属流水级的配置，从而在保持 Megatron 单卡前向/反向语义不变的前提下，实现按流水级异构的重计算。

\paragraph{与 MMoH 的协同。}
HR-MMoH 以 MMoH 给出的流水级划分与设备映射为输入，二者形成层次化协作：若在统一检查点策略下某流水级仍出现显存溢出，可将该信息反馈至 MMoH 规划器（例如在显存约束中提高该流水级的权重，或对该流水级进行更细粒度的块划分），从而实现划分与显存配置的迭代优化。第三章所述的提前终止调度器在冻结流水级上省略反向计算与梯度通信，与各流水级独立的检查点配置无冲突：冻结流水级可采用较保守的检查点策略以节省计算，非冻结流水级则按显存与算力约束进行均衡配置。

\section{实验验证}
\label{sec:HR-MMoH-exp}

本节从实验设置、实验结果与实验分析三方面对 HR-MMoH 的设计动机与配置思路进行说明；完整端到端实现及与 MMoH 的联合吞吐对比可作为后续工作深化。

\subsection{实验设置}

\paragraph{平台与配置思路验证。}
实验环境与第三章 MMoH 实验一致，采用两种异构集群（Cluster-S、Cluster-L）及两种 MLLM（MLLM-3B、MLLM-12B）。为验证“按流水级差异化配置”的合理性，对 MMoH 划分得到的各流水级进行显存与耗时剖析：在每流水级对应的设备上，分别测量 \texttt{recompute\_granularity=None}、\texttt{selective} 且 \texttt{recompute\_modules=[core\_attn]}、\texttt{[core\_attn, mlp]}、\texttt{[core\_attn, mlp, layernorm]} 以及 \texttt{full} + \texttt{uniform} 且不同 \texttt{recompute\_num\_layers} 下的激活峰值与单流水级前向+反向时间（含重算）。据此可得到各流水级在满足显存约束下的可选配置集合及对应的 $T_s$，用于验证显存紧张流水级与充裕流水级在“最省显存”与“最少重算”配置上的差异是否显著，以及按流水级赋配置后各 $T_s$ 是否更均衡。

\paragraph{对比与指标。}
重计算的对比基线为全局统一配置：所有流水级不采用重计算，记为\texttt{no-recomp}；所有流水级采用相同的 \texttt{full recomputation}，即在前向传播中丢弃所有的中激活值，记为\texttt{full}；
所有流水级采用相同的 \texttt{selective} + \texttt{[core\_attn]} 或 \texttt{[core\_attn, mlp, layernorm]}（记为 \texttt{select-1} 与 \texttt{select-2}）。
指标包括：（1）各流水级在满足显存不溢出条件下的显存占用峰值与单流水级耗时 $T_s$；（2）完整训练循环的训练吞吐（\texttt{samples/s} 和迭代时间\texttt{s/iteration}。

\subsection{实验结果}

在 MMoH 对 MLLM-3B / MLLM-12B 的划分结果上，对各流水级进行上述剖析后可观察到：显存相对紧张流水级（如承载 LLM 中段或大 block 的流水级）在 \texttt{recompute\_granularity=None} 或仅 \texttt{[core\_attn]} 时易接近或超出设备显存上限，需采用 \texttt{[core\_attn, mlp, layernorm]} 或 \texttt{full} + \texttt{uniform} 方能安全训练；而显存充裕流水级（如部分编码器或投影层所在流水级）采用 \texttt{[core\_attn, mlp, layernorm]} 会带来不必要的重算，单流水级耗时 $T_s$ 明显增加。按流水级赋予差异化配置（显存紧张流水级用较激进检查点，充裕流水级用较保守检查点）后，各流水级 $T_s$ 的方差较 Uniform-1 / Uniform-3 有所下降，迭代时间 $\max_s T_s$ 在多数场景下优于或接近“为照顾最紧张流水级而全局采用 Uniform-3”的配置，说明异构感知的配置策略在理论上能够兼顾显存安全与负载均衡。完整训练循环下的吞吐对比有待在 HR-MMoH 与 Megatron/MMoH 的深度集成实现后进一步报告。

\subsection{实验分析}

\paragraph{显存—算力权衡。}
异构集群上各流水级显存与算力分布不均，统一配置要么为防 OOM 而整体偏保守导致重算过多，要么为减重算而整体偏激进导致瓶颈流水级 OOM。按流水级配置后，可在满足各流水级显存约束的前提下，将“重算预算”更多分配给显存紧张、算力相对充裕的流水级，从而在系统层面降低迭代时间。

\paragraph{与 MMoH 划分的配合。}
MMoH 的复合粒度与冻结感知划分已使各流水级负载趋于均衡；HR-MMoH 在此基础上对显存维度做细粒度调节，二者结合可形成“划分—调度—显存优化”的闭环。未来工作可将 HR-MMoH 的配置搜索与 MMoH 的整数规划联合建模，或采用分层决策：先由 MMoH 确定划分与映射，再由 HR-MMoH 模块为各流水级输出推荐配置并可选地反馈显存不足信息以触发重新划分。

\section{本章小结}
\label{sec:HR-MMoH-summary}

本章在 MMoH 的异构流水线划分与提前终止调度基础上，针对显存受限流水级仍可能 OOM、而现有重计算策略多采用全局统一配置难以兼顾显存与效率的问题，
介绍了 HR-MMoH 的异构重计算思路。
首先系统梳理了 Megatron-LM 的重计算配置参数（\texttt{recompute\_granularity}、\texttt{recompute\_method}、\texttt{recompute\_num\_layers}、\texttt{recompute\_modules} 等）
及其 \texttt{full/selective} 粒度与层级划分方式，并指出其在同构假设下的局限；随后给出了 HR-MMoH 的设计要点：以各流水级显存上限、算力为约束，为各流水级独立配置重计算粒度与子模块集合，通过基于剖析的配置策略在满足显存不溢出的前提下最小化迭代时间并平衡各流水级负载，并与 MMoH 划分及提前终止调度器协同。
实验部分说明了配置思路的验证方式与预期效果，并指出完整端到端实现与联合吞吐对比可作为后续工作。
HR-MMoH 与第三章的 MMoH 共同构成“划分—调度—显存优化”的联合优化框架，为异构 GPU 集群上 MLLM 的高效训练提供了显存维度的补充手段。
